\subsection[Category of Templates]{The category of templates $\tempC$}

The first step towards introducing template heaps requires us to distinguish between two different kinds of locations.
We have (i) concrete locations that are associated with some cell contents and 
(ii) location variables that stand for \emph{any arbitrary cell contents}.

To model these two kinds of locations we make use of the category of templates defined as the arrow category $\arrC\worldC$
\begin{itemize}
    \item Objects: Sort-preserving injective renamings between worlds $\Minus\vw \xto\sigma \Plus\vw$.
    \item Morphisms: Injective, sort-preserving $\concL\rho : \Minus\vw \to \Plus{\vw'}$ and $\anyL\rho : \Minus\vw \to \Plus{\vw'}$, such that: 
                    \begin{center}
                    \begin{tikzcd}{\Plus\vw} \ar[r, tail, "{\anyL\rho}"] & {\Plus{\vw'}} \\
                                  {\Minus\vw} \ar[u, tail, "\sigma"] \ar[r, tail, "{\concL\rho}"'] & {\Minus{\vw'}} \ar[u, tail, "\sigma'"']
                     \end{tikzcd} 
                    \end{center}
\end{itemize}
We have two forgetful functors $\concL\dash, \anyL\dash : \tempC \to \worldC$, projecting the $-$ and $+$ components respectively.

\begin{notation}
\noindent
\begin{itemize}
    \item We refer to the elements of $\concL\sigma$ as the concrete locations of $\sigma \in \objC{\tempC}$.
    \item We refer to the elements of $\Plus\vw \ominus \sigma$ as the location variables of $\sigma \in \objC{\tempC}$, notation $\varL\sigma$.
    \item We refer to elements of $\anyL\sigma$ as \emph{all} the locations of $\sigma \in \objC{\tempC}$.
    \item If it is clear from the context we overload the notations $\concL\sigma = \concL\sigma \xto\sigma \anyL\sigma$ and
        $\varL\sigma = \varL\sigma \xto{\compl\sigma} \anyL\sigma$ for the inclusions into $\anyL\sigma$.
    \item We have an induced morphism $\varL\rho : \varL\sigma \xto{\varL\sigma} \anyL\sigma \xto{\anyL\rho} \anyL{\sigma'}$.
    \item We refer to elements $\anyL{\sigma'} \ominus \anyL\rho$ as the new locations introduced by $\rho$, 
          and $\concL{\sigma'} \ominus \concL\rho$ as the concrete locations introduced by $\rho$.
    \item We refer to the variables $\varL\sigma$ that are instantiated by $\rho$ as $\concL{\sigma'} \pltimes{\anyL{\sigma'}} \varL\sigma$: 
        \begin{center}
            \begin{tikzcd}
                {\concL{\sigma'} \pltimes{\anyL{\sigma'}} \varL\sigma} \ar[r] \ar[d] \pbsquare{dr} & {\varL\sigma} \ar[d, "\varL\rho"] \\
                {\concL{\sigma'}} \ar[r, "\concL{\sigma'}"']                    & {\anyL{\sigma'}}
            \end{tikzcd}
        \end{center}
            which are exactly the $\varL\sigma$ that are mapped to concrete locations $\concL{\sigma'}$ by $\rho$.
    \item We refer to the variables $\varL{\sigma'}$ that are fresh w.r.t. $\rho$ as $\varL{\sigma'} \pltimes{\anyL{\sigma'}} (\anyL{\sigma'} \ominus \anyL\rho)$: 
        \begin{center}
            \begin{tikzcd}
                {\varL{\sigma'} \pltimes{\anyL{\sigma'}} (\anyL{\sigma'} \ominus \anyL\rho)}
                    \ar[r] \ar[d] \pbsquare{dr} 
                        & {\anyL{\sigma'} \ominus \anyL\rho} 
                            \ar[d, "\compl{\anyL\rho}"] 
                \\
                {\varL{\sigma'}}    
                    \ar[r, "\varL{\sigma'}"'] 
                        & {\anyL{\sigma'}} 
            \end{tikzcd}
        \end{center}
            which are exactly the $\varL{\sigma'}$ that are not mapped from any location $\anyL\sigma$ by $\rho$.
    \item Similarly, $\varL{\sigma'} \pltimes{\anyL{\sigma'}} \varL\sigma$ refers to the variables $\varL{\sigma'}$ that are mapped from variables in $\varL\sigma$ by $\rho$. 
        \begin{center}
            \begin{tikzcd}
                {\varL{\sigma'} \pltimes{\anyL{\sigma'}} \varL\sigma} \ar[r] \ar[d] \pbsquare{dr} &  {\varL\sigma} \ar[d, "\varL\rho"] \\
                {\varL{\sigma'}} \ar[r, "\varL{\sigma'}"'] & {\anyL{\sigma'}} 
            \end{tikzcd}
        \end{center}
    \item Notice that from $\concL{\sigma'}$ and $\varL{\sigma'}$ being disjoint, we have that concrete locations in $\concL\sigma$ cannot become uninstantiated
        \begin{center}
            \begin{tikzcd}
                {\emptyset} \ar[rr] \ar[d] \pbsquare{drr} && {\varL{\sigma'}} \ar[d, "\varL{\sigma'}"] \\
                {\concL\sigma} \ar[r, "\concL\rho"'] & {\concL{\sigma'}} \ar[r, "\concL{\sigma'}"'] & {\anyL{\sigma'}} 
            \end{tikzcd}
        \end{center}
    \end{itemize}
\end{notation}

\begin{example}\label{ex:tempC-overview}
    Consider the following morphism in $\tempC$
    \begin{center}
    \begin{tikzcd}
        {\va : \vC, \vb : \vC'} \ar[rr, hook, "\anyL\rho = \incl^{\va,\vb \inclto \va,\vb,\vc}" yshift=5pt]               && {\va : \vC, \vb : \vC', \vc : \vC''} \\
        {\va : \vC} \ar[u, hook, "\incl^{\va \inclto \va,\vb}"] \ar[rr, hook, "\concL\rho"'] && {\va : \vC, \vb : \vC'} \ar[u, hook, "\incl^{\va,\vb \inclto \va,\vb,\vc}"']
    \end{tikzcd}
    \end{center}
    \begin{itemize}
        \item We have variables $\varL{\incl^{\va \inclto \va,\vb}} = \{ \vb : \vC' \}$, and $\varL{\incl^{\va,\vb \inclto \va,\vb,\vc}} = \{ \vc : \vC''\}$.
        \item Variables instantiated by $\rho$: 
                \[\concL{\incl^{\va,\vb \inclto \va,\vb,\vc}} \pltimes{\{ \va, \vb, \vc \}} \varL{\incl^{\va \inclto \va,\vb}} = \{ \vb : \vC' \} \]
        \item Variables $\varL{\incl^{\va,\vb \inclto \va,\vb,\vc}}$ that are fresh w.r.t. to $\rho$: 
                \[ \varL{\incl^{\va,\vb \inclto \va,\vb,\vc}} \pltimes{\{ \va,\vb,\vc\}} (\anyL{\incl^{\va,\vb \inclto \va,\vb,\vc}} \ominus \anyL\rho) = \{ \vc : \vC''\} \]
        \item Variables $\varL{\incl^{\va,\vb \inclto \va,\vb,\vc}}$ mapped from variables in $\varL{\incl^{\va \inclto \va,\vb}}$ by $\rho$: 
                \[ \varL{\incl^{\va,\vb \inclto \va,\vb,\vc}} \pltimes{\{ \va,\vb,\vc\}} \varL{\incl^{\va \inclto \va,\vb}} = \emptyset \]
    \end{itemize}
\end{example}

\begin{definition}[Variable Extension Functor]
We define the variable extension functor as a means to add $\vw$-many variables.
\[\begin{array}{r@{\,}c@{\,}l}
    \extF\dash\vw & : & \tempC \to \tempC \\
    \extF{(\concL\sigma \xto{\concL\sigma} \anyL\sigma)}{\vw} & = & \begin{tikzcd}{\anyL\sigma \mplus \vw} \\ 
                                                                                  {\concL\sigma} \ar[u, "\incl_1 \comp \concL\sigma"']
                                                                    \end{tikzcd} \\
    \extF\rho\vw & = & \pair{\concL\me}{\anyL\me \mplus 1_\vw}
\end{array}\]
Further, for every functor $\vG : \tempC \to \setC$ and $\vw$, we have a $\extNat\vG\vw : \vG \To \vG\extF\dash{\vw'}$.
\end{definition}
\begin{example}
    The action of $\extF\dash{\{ \vd : \vC''' \}}$ on morphism $\rho$ from \cref{ex:tempC-overview} is $\extF\rho{\{ \vd : \vC'''\}}$ defined as
    \begin{center}
    \begin{tikzcd}
        {\{ \va : \vC, \vb : \vC' \} \oplus \{ \vd : \vC''' \} } 
            \ar[r, "\anyL\rho \oplus 1_\vd" yshift=2pt] 
            & {\{ \va : \vC, \vb : \vC', \vc : \vC'' \} \oplus \{ \vd : \vC'''\}} 
        \\
        {\va : \vC} 
            \ar[u, hook, "\incl_1 \comp \incl^{\va \inclto \va,\vb}"] \ar[r, hook, "\concL\rho"'] 
            & {\va : \vC, \vb : \vC'} \ar[u, hook, "\incl_1 \comp \incl^{\va,\vb \inclto \va,\vb,\vc}"']
    \end{tikzcd}
    \end{center}
\end{example}

\begin{lemma}[Monomorphisms in $\tempC$]\label{lem:monos-T}
    Every morphism in $\tempC$ is mono.
\end{lemma}
\begin{proof}Follows directly from injectivity of underlying morphisms.\end{proof}

\begin{lemma}[Initial Object in $\tempC$]
    The category $\tempC$ has an initial object.
\end{lemma}
\begin{proof}Consider template $?_\emptyset : \emptyset \to \emptyset$.\end{proof}

\begin{lemma}[Equalizers in $\tempC$]
    The category $\tempC$ has all equalizers, and $\concL\dash$, $\anyL\dash$ preserve equalizers.
\end{lemma}
\begin{proof}
    The equalizer
    \begin{tikzcd}{\eqz{\rho_1}{\rho_2}} \ar[r] & {\sigma_1} \ar[r, yshift=3pt, "\rho_1"] \ar[r, yshift=-3pt, "\rho_2"'] & {\sigma_2}\end{tikzcd} 
    is given componentwise
    \begin{center}
    \begin{tikzcd}
        {\concL{\hat\sigma}} 
            \ar[r, dashed, "\concL\vh"] \ar[rrr, yshift=4pt, bend left=30, "\concL{\hat\rho}"] \ar[d, "\concL{\hat\sigma}"]
            & {\eqz{\concL{\rho_1}}{\concL{\rho_2}}} 
                \ar[rr, "\eqz{\concL{\rho_1}}{\concL{\rho_2}}"] \ar[d, dashed] 
                    && {\concL{\sigma_1}} \ar[d, "\concL{\sigma_1}"]
        \\
        {\anyL{\hat\sigma}}
            \ar[r, dashed, "\anyL\vh"] \ar[rrr, bend right=30, "\anyL{\hat\rho}"]
            & {\eqz{\anyL{\rho_1}}{\anyL{\rho_2}}} 
                \ar[rr, "\eqz{\anyL{\rho_1}}{\anyL{\rho_2}}"]
                    && {\anyL{\sigma_1}} \ar[r, yshift=3pt, "\anyL{\rho_1}"] 
                        \ar[r, yshift=-3pt, "\anyL{\rho_2}"'] 
                            & {\anyL{\sigma_2}}
    \end{tikzcd} 
    \end{center}
\end{proof}
\begin{example}
    This is an equalizer diagram 
    \begin{center}
    \begin{tikzcd}
        {?_{\vb}} \ar[r, "\pair{?_\va}{\incl_2}"] 
            & {\incl_1^{\va \to \va \oplus \vb}}
                \ar[rr, yshift=3pt, "\pair{\incl_1}{\incl_1 \oplus 1_\vb}"] \ar[rr, yshift=-3pt, "\pair{\incl_2}{\incl_2 \oplus 1_\vb}"'] 
            && {\incl_1^{\va_1 \oplus \va_2 \to \va_1 \oplus \va_2 \oplus \vb}}
    \end{tikzcd}
    \end{center}
\end{example}

\begin{lemma}[Pullbacks in $\tempC$]\label{lem:pullbacks-tempC}
    The category $\tempC$ has all pullbacks and $\concL\dash$, $\anyL\dash$ preserve pullbacks.
\end{lemma}
\begin{proof}
    The pullback of $\sigma_1 \xto{\rho_1} \hat\sigma \xfrom{\rho_2} \sigma_2$ is given componentwise
    \begin{center}
    \begin{tikzcd}
        {\concL\sigma}
            \ar[rrrr, bend left=30, "\concL{\rho_2'}"]
            \ar[ddrr, bend right=30, "\concL{\rho_1'}"']
            \ar[dr, "\concL{\hat\sigma}"]
            \ar[rr, dashed, "\concL\vh"]
                && {\concL{\sigma_1} \pltimes{\concL\sigma} \concL{\sigma_2}}
                \ar[rr, "\ext^{\concL{\rho_1}}{\concL{\rho_2}}"]
                    \ar[dr, dashed]
                    \ar[dd, "\ext^{\concL{\rho_2}}\concL{\rho_1}"' yshift=20pt]
                        && {\concL{\sigma_2}}
                            \ar[dr, "\concL{\sigma_2}"']
                            \ar[dd, "\concL\rho"' yshift=20pt]
        \\
            & {\anyL\sigma}
                \ar[rrrr, bend left=60, "\anyL{\rho_2'}" xshift=20pt]
                \ar[rr, dashed, "\anyL\vh" xshift=-40pt]
                \ar[ddrr, bend right=30, "\anyL{\rho_1'}"']
                    && {\anyL{\sigma_1} \pltimes{\anyL\sigma} \anyL{\sigma_2}}
                        \ar[rr, "\ext^{\anyL{\rho_2}}\anyL{\rho_1}" xshift=-20pt]
                        \ar[dd, "\ext^{\anyL{\rho_1}}\anyL{\rho_2}" yshift=20pt]
                                && {\anyL{\sigma_2}}
                                    \ar[dd, "\anyL{\rho_2}"]
        \\
                && {\concL{\sigma_1}}
                    \ar[rr, "\concL{\rho_1}" xshift=-40pt]
                    \ar[rd, "\ext^{\anyL{\rho_1}}\anyL{\rho_2}"']
                        && {\concL\sigma}
                            \ar[dr, "\concL\sigma"]
        \\
                    &&& {\anyL{\sigma_1}} 
                            \ar[rr, "\anyL{\rho_1}"]
                            && {\anyL\sigma}
    \end{tikzcd}
    \end{center}
\end{proof}

\begin{example}
    This is a pullback square
    \begin{center}
    \begin{tikzcd} 
        {?_\vb} 
            \ar[r, "\pair{?_{\va_1}}{\incl_2}"] 
            \ar[d, "\pair{?_{\va_2}}{\incl_2}"']
            \pbsquare{dr}
            & {\incl_1^{\va_1 \to \va_1 \oplus \vb}}
                \ar[d, "\pair{\incl_1}{1_{\va_1} \oplus \incl_2}"]
        \\
        {1_{\va_2 \oplus \vb}}
            \ar[r, "\pair{\incl_2}{\incl_2}"']
            & {1_{\va_1 \oplus (\va_2 \oplus \vb)}}
    \end{tikzcd}
    \end{center}
\end{example}

\subsubsection{Generated Cospans}

\begin{lemma}[Generated Cospans in $\tempC$]\label{lem:gen-cospan}
    Given cospan $\anyL{\sigma_1} \xto{\tau_1} \hat\vw \xfrom{\tau_2} \anyL{\sigma_2}$ in $\worldC$ 
    we generate a cospan $\sigma_1 \xto{\rho^\tau_1} \sigma^\tau \xfrom{\rho^\tau_2} \sigma_2$ in the least possible way, 
    that is, it satisfies
    \begin{itemize}
        \item Satisfy $\anyL{\sigma^\tau} = \hat\vw$, $\anyL{\rho^\tau_1} = \tau_1$ and $\anyL{\rho^\tau_2} = \tau_2$.
        \item For any other cospan $\sigma_1 \xto{\rho_1'} \hat\sigma \xfrom{\rho_2'} \sigma_2$ that factors through $\hat\vw$
                \begin{tikzcd}
                    {\anyL{\sigma_1}} \ar[dr, "\anyL{\rho_1'}"'] \ar[r, "\tau_1"] & {\hat\vw} \ar[d, dashed, "\exists \tau'"] & {\anyL{\sigma_2}} \ar[l, "\tau_2"'] \ar[dl, "\anyL{\rho_2'}"] \\
                                                                                      & {\anyL{\hat\sigma}}
                \end{tikzcd},
              there exists a unique $\vh : \sigma^\tau \to \hat\sigma$ that satisfies $\rho_1' = \vh \comp \rho^\tau_1$ and $\rho_2' = \vh \comp \rho^\tau_2$.
    \end{itemize}

    We say that $\rho^\tau_1, \rho^\tau_2$ are generated by $\tau_1, \tau_2$.
    Intuitively, a generated cospan selects the smallest $\concL{\sigma^\tau}$ possible, that is, it selects the least concrete locations possible.
\end{lemma}
\begin{proof}
    We define $\anyL{\sigma^\tau} = \hat\vw$, $\anyL{\rho^\tau_1} = \tau_1$, $\anyL{\rho^\tau_2} = \tau_2$.
    To define $\concL{\sigma^\tau}$ we first consider their common concrete locations through the pullback
    \begin{center}
    \begin{tikzcd}[sep=small]
        {\concL{\sigma_1} \pltimes{\hat\vw} \concL{\sigma_2}} 
            \ar[rrrr, "\res_{\concL{\sigma_2};\tau_2}\concL{\sigma_1};\tau_1"] 
            \ar[dd, "\res_{\concL{\sigma_1}; \tau_1}\concL{\sigma_2}; \tau_2"'] 
            \pbsquare{ddrrrr}
                &&&& {\concL{\sigma_2}} \ar[d, "\concL{\sigma_2}"] \\
                &&&& {\anyL{\sigma_2}} \ar[d, "{\tau_2}"] \\
        {\concL{\sigma_1}} \ar[r, "\concL{\sigma_1}"'] 
                & {\anyL{\sigma_1}} \ar[rrr, "{\tau_1}"']
                &&& {\hat\vw}
    \end{tikzcd}
    \end{center}
    and further take the L.I.C. in $\worldC$ of the pullback morphisms
    \begin{center}
    \begin{tikzcd}
        {\concL{\sigma_1} \pltimes{\hat\vw} \concL{\sigma_2}} 
            \ar[rr, "\res_{\concL{\sigma_2}; \tau_2}\concL{\sigma_1}; \tau_1" yshift=5pt] 
            \ar[dd, "\res_{\concL{\sigma_1}; \tau_1}\concL{\sigma_2}; \tau_2"']
                && {\concL{\sigma_2}} \ar[dd, "\concL{\rho^\tau_2}"] 
    \\\\
        {\concL{\sigma_1}} \ar[rr, "\concL{\rho^\tau_1}"']
                && {\concL{\sigma^\tau} := \concL{\sigma_1} \liplus{\concL{\sigma_1} \pltimes{\hat\vw} \concL{\sigma_2}} \concL{\sigma_2}}
                                           
    \end{tikzcd}
    \end{center}
    The morphism $\concL{\sigma^\tau} \xto{\concL{\sigma^\tau}} \anyL{\sigma^\tau}$ is the unique morphism induced by the L.I.C. towards the pullback. 
    The UP gives that the morphisms $\rho^\tau_1$, $\rho^\tau_2$ are commuting squares.

    \paragraph{Least Property}
    Consider a compatible cospan $\rho_1', \rho_2'$.
    We only need to extend $\anyL\vh = \tau'$ to a morphism in $\tempC$.
    Let $\concL\vh$ be the induced morphism by the L.I.C.
    \begin{center}
    \begin{tikzcd}
        {\concL{\sigma_1} \pltimes{\hat\vw} \concL{\sigma_2}} \ar[d, "\res_{\concL{\sigma_1}; \tau_1}\concL{\sigma_2}; \tau_2"']
                                                                              \ar[rr, "\res_{\concL{\sigma_2}; \tau_2}\concL{\sigma_1}; \tau_1"]
            && {\concL{\sigma_2}} \ar[d, "\concL{\rho^\tau_2}"] \ar[rdd, bend left=40, "\concL{\rho_2'}"] & \\
        {\concL{\sigma_1}} \ar[rr, "\concL{\rho^\tau_1}"'] \ar[rrrd, bend right=20, "\concL{\rho_1'}"'] 
            && {\concL{\sigma^\tau}} \ar[dr, dashed, "\concL\vh"] & \\
                                    &&& {\concL{\hat\sigma}}
    \end{tikzcd}
    \end{center}
    The outer square commutes by
    \begin{align*}
        \concL{\rho_2'} \comp (\res_{\concL{\sigma_2};\tau_2}\concL{\sigma_1};\tau_1) & = \concL{\rho_1'} \comp (\res_{\concL{\sigma_1};\tau_1}\concL{\sigma_2};\tau_2) \iff \\
        \concL{\hat\sigma} \comp \concL{\rho_2'} \comp (\res_{\concL{\sigma_2};\tau_2}\concL{\sigma_1};\tau_1) & = \concL{\hat\sigma} \comp \concL{\rho_1'} \comp (\res_{\concL{\sigma_1};\tau_1}\concL{\sigma_2};\tau_2) = \\
        \anyL{\rho_2'} \comp \concL{\sigma_2} \comp (\res_{\concL{\sigma_2};\tau_2}\concL{\sigma_1};\tau_1) & = \anyL{\rho_2'} \comp \concL{\sigma_1} \comp (\res_{\concL{\sigma_1};\tau_1}\concL{\sigma_2};\tau_2) = \\
        \tau' \comp \tau_2 \comp \concL{\sigma_2} \comp (\res_{\concL{\sigma_2};\tau_2}\concL{\sigma_1};\tau_1) & = \tau' \comp \tau_1 \comp \concL{\sigma_1} \comp (\res_{\concL{\sigma_1};\tau_1}\concL{\sigma_2};\tau_2)
    \end{align*}
    and satisfies independence by \cref{lem:pullback-mono-post-comp}.
    Further $\vh$ is a commuting square, that is $\concL{\hat\sigma} \comp \concL\vh = \anyL\vh \comp \concL{\sigma^\tau}$, by the jointly epi property of $\concL{\rho^\tau_1}$ and $\concL{\rho^\tau_2}$.
\end{proof}

\begin{lemma}[Generated Cospans and Jointly Epi Property]\label{lem:gen-cospan-epi-prop}
    A cospan $\sigma_1 \xto{\rho_1} \hat\sigma \xfrom{\rho_2} \sigma_2$ is 
    generated by $\anyL{\sigma_1} \xto{\anyL{\rho_1}} \anyL{\hat\sigma} \xfrom{\anyL{\rho_2}} \anyL{\sigma_2}$ 
    iff $\concL{\rho_1}$ and $\concL{\rho_2}$ are jointly surjective.
\end{lemma}
\begin{proof}\nl
    $\to$) 
    Generated cospans are unique up to isomorphism, and so $\rho_1$, $\rho_2$ is iso to the cospan induced by \cref{lem:gen-cospan} at $\anyL{\rho_1}$, $\anyL{\rho_2}$. 
    Since the cospan of \cref{lem:gen-cospan} at $\concL\dash$ is a L.I.C., we have that $\concL{\rho_1}$, $\concL{\rho_2}$ are jointly epi.

    $\from$)
    We only need to show that the cospan $\rho_1$, $\rho_2$ satisfies the UP of generated cospans, so consider
    \begin{center}
    \begin{tikzcd}
        {\anyL{\sigma_1}} 
            \ar[dr, "\anyL{\rho_1'}"'] \ar[r, "\anyL{\rho_1}"] 
            & {\anyL{\hat\sigma}} 
                \ar[d, dashed, "\exists \anyL\vh"] 
                & {\anyL{\sigma_2}} 
                    \ar[l, "\anyL{\rho_2}"'] \ar[dl, "\anyL{\rho_2'}"] 
        \\
            & {\anyL{\sigma'}}
    \end{tikzcd}
    \end{center}
    Define $\concL\vh$ as 
    \begin{align*}
        \concL\vh & : \concL{\hat\sigma} \to \concL{\sigma'} \\
        \concL\vh(\va : \vC) & \eqdef \begin{cases} \concL{\rho_1'}(\va_1)  & \;\mif\; \concL{\rho_1}(\va_1) = \va \\
                                                    \concL{\rho_2'}(\va_2)  & \;\mif\; \concL{\rho_2}(\va_2) = \va
                                      \end{cases}
    \end{align*}
    It is well-defined for if $\concL{\rho_1}(\va_1) = \va = \concL{\rho_1}(\va_2)$, we have
    \begin{align*}
        \concL{\hat\sigma}\concL{\rho_1}(\va_1) &= \concL{\hat\sigma}\concL{\rho_2}(\va_2) \iff \\
        \anyL{\rho_1}\concL{\sigma_1}(\va_1) &= \anyL{\rho_2}\concL{\sigma_2}(\va_2) \iff \\
        \anyL\vh\anyL{\rho_1}\concL{\sigma_1}(\va_1) &= \anyL\vh\anyL{\rho_2}\concL{\sigma_2}(\va_2) \iff \\
        \anyL{\rho_1'}\concL{\sigma_1}(\va_1) &= \anyL{\rho_2'}\concL{\sigma_2}(\va_2) \iff \\
        \concL{\sigma'}\concL{\rho_1'}(\va_1) &= \concL{\sigma'}\concL{\rho_1'}(\va_2) \iff \\
        \concL{\rho_1'}(\va_1) &= \concL{\rho_1'}(\va_2)
    \end{align*}
    It is easy to show that $\vh$ is a morphism in $\tempC$ and uniqueness of $\vh$ follows from the jointly epi property. 
\end{proof}

\begin{lemma}[Instantiation data for Generated Cospans]\label{lem:gen-cospan-data}
    Given a diagram of the form 
        \begin{tikzcd} 
            & {\anyL\sigma} \ar[dl, "\anyL{\rho_1}"'] \ar[dr, "\anyL{\rho_2}"] & \\ 
            {\anyL{\sigma_1}} \ar[dr, "\tau_1"'] && {\anyL{\sigma_2}} \ar[dl, "\tau_2"] \\
                                                & {\hat\vw} &
        \end{tikzcd} where $\rho_2$ has data $\mvarM : \dheapF(\concL{\sigma_2} \ominus \concL{\rho_2}, \anyL{\sigma_2})$,
        we can equip $\rho^\tau_1$ with data $\eta_{\rho^\tau_1} : \dheapF(\concL{\hat\sigma} \ominus \concL{\rho^\tau_1}, \anyL{\hat\sigma})$.
\end{lemma}
\begin{proof}
    We have a canonical isomorphism
    \[ \concL{\hat\sigma} \ominus \concL{\rho^\tau_1} \iso \concL{\sigma_2} \ominus \res_{(\concL{\sigma_2}; \tau_2)}\concL{\sigma_1};\tau_1 \]
    and further from the pullback property
    \begin{center}
    \begin{tikzcd}[sep=small]
    {\concL\sigma} \ar[rrrrrd, bend left=20, "\concL{\rho_2}"] \ar[dddr, bend right=50, "\concL{\rho_1}"'] \ar[dr, dashed, "\vh"] & 
\\
        & {\concL{\sigma_1} \pltimes{\hat\vw} \concL{\sigma_2}} 
            \ar[rrrr, "\res_{\concL{\sigma_2};\tau_2}\concL{\sigma_1};\tau_1"] 
            \ar[dd, "\res_{\concL{\sigma_1}; \tau_1}\concL{\sigma_2}; \tau_2"] 
            \pbsquare{ddrrrr}
                &&&& {\concL{\sigma_2}} \ar[d, "\concL{\sigma_2}"] 
\\
                &&&&& {\anyL{\sigma_2}} \ar[d, "{\tau_2}"] 
\\
        & {\concL{\sigma_1}} \ar[r, "\concL{\sigma_1}"'] 
            & {\anyL{\sigma_1}} \ar[rrr, "{\tau_1}"']
                &&& {\hat\vw}
    \end{tikzcd}
    \end{center}
    we have 
\[ \concL{\sigma_2} \ominus \concL{\rho_2} \iso (\concL{\sigma_2} \ominus \res_{(\concL{\sigma_2}; \tau_2)}\concL{\sigma_1};\tau_1) \mplus 
                                                (\concL{\sigma_1} \pltimes{\hat\vw} \concL{\sigma_2} \ominus \vh) \]
    Putting these together gives a morphism 
    \[ \concL{\hat\sigma} \ominus \concL{\rho^\tau_1} \xto{\incl_1} \concL{\sigma_2} \ominus \concL{\rho_2} \]
    and so we define $\mvarM_{\rho^\tau_1} = \dheapF(\incl_1, \tau_2)(\mvarM)$.
\end{proof}

\subsubsection{Local Independence Structure}

We equip $\tempC$ with a local independence structure that takes into consideration the instantiations of location variables.

\begin{definition}[Local Independent Squares in $\tempC$]
    \begin{tikzcd}[sep=small] 
                    & {\sigma} \ar[ld, "\rho_1"'] \ar[rd, "\rho_2"] & \\
        {\sigma_1} \ar[rd, "\rho_1'"'] &            & {\sigma_2} \ar[ld, "\rho_2'"] \\
                                &  {\hat\sigma} &
    \end{tikzcd} 
    is a local independent square when the pullback of $\anyL{\sigma_1} \pltimes{\anyL{\hat\sigma}} \anyL{\sigma_2}$ is $\anyL\sigma$
    \begin{center}
        \begin{tikzcd}
            {\anyL\sigma} \ar[r, "\anyL{\rho_2}"] \ar[d, "\anyL{\rho_1}"'] \pbsquare{dr} & {\anyL{\sigma_2}} \ar[d, "\anyL{\rho'_2}"] \\
            {\anyL{\sigma_1}} \ar[r, "\anyL{\rho'_1}"']                                  & {\anyL{\hat\sigma}}
        \end{tikzcd}
    \end{center}
    Thus an independent square conflates two locations $\va \in \anyL{\sigma_1}$ and $\vb \in \anyL{\sigma_2}$ only if they come from $\anyL\sigma$.
\end{definition}

\begin{lemma}[Local Independent Coproducts in $\tempC$]\label{lem:lic-exist}
    For every span $\sigma_1 \xfrom{\rho_1} \sigma \xto{\rho_2} \sigma_2$ the local independent coproduct (L.I.C.) exists.
\end{lemma}
\begin{proof}
    Consider the L.I.C. in $\worldC$
    \begin{center}
    \begin{tikzcd}
        {\anyL\sigma} \ar[d, "{\anyL{\rho_2}}"'] \ar[rr, "{\anyL{\rho_1}}"] && {\anyL{\sigma_1}} \ar[d, "{\ext^{\anyL{\rho_1}}\anyL{\rho_2}}"] \\
        {\anyL{\sigma_2}} \ar[rr, "{\ext^{\anyL{\rho_2}}\anyL{\rho_1}}"'] && {\anyL{\sigma_1} \liplus{\anyL\sigma} \anyL{\sigma_2}}
    \end{tikzcd}
    \end{center}

    Using \cref{lem:gen-cospan} for $\anyL{\sigma_1} \xto{\ext^{\anyL{\rho_1}}\anyL{\rho_2}} \anyL{\sigma_1} \liplus{\anyL\sigma} \anyL{\sigma_2} \xfrom{\ext^{\anyL{\rho_2}}\anyL{\rho_1}} \anyL{\sigma_2}$ can extend this to morphisms in $\tempC$
        \[ \sigma_1 \xto{\ext^{\rho_1}\rho_2} \sigma_1 \liplus\sigma \sigma_2 \xfrom{\ext^{\rho_2}\rho_1} \sigma_2 \]
        We have that $\ext^{\rho_1}\rho_2 \comp \rho_1 = \ext^{\rho_2}\rho_1 \comp \rho_2$ from $\ext^{\anyL{\rho_1}}\anyL{\rho_2} \comp \anyL{\rho_1} = \ext^{\anyL{\rho_2}}\anyL{\rho_1} \comp \anyL{\rho_2}$.

    \paragraph{Induced Morphism}
    Consider the diagram where each square is independent.
    \begin{center}
    \begin{tikzcd}     
                                          & {\sigma} \ar[ld, "{\rho_1}"'] \ar[rd, "{\rho_2}"] & \\
{\sigma_1} \ar[rd, pos=0.30, "{\ext^{\rho_1}\rho_2}"] \ar[rdd, "\rho_1'"'] &   & {\sigma_2} \ar[ld, pos=0.30,"{\ext^{\rho_2}\rho_1}"'] \ar[ldd, "\rho_2'"] \\
                                                             & {\sigma_1 \liplus\sigma \sigma_2} \ar[d, dashed, pos=0.20, "{\hat\rho}"] & \\
                                          & {\hat\sigma} &
    \end{tikzcd}
    \end{center}

    The $\anyL{\hat\rho}$ part of morphism $\hat\rho$ is induced by
    \begin{center}
    \begin{tikzcd}
        {\anyL\sigma} \ar[d, "{\anyL{\rho_2}}"'] \ar[rr, "{\anyL{\rho_1}}"] 
            && {\anyL{\sigma_1}} \ar[d, "{\ext^{\anyL{\rho_1}}\anyL{\rho_2}}"] \ar[rdd, bend left=40, "\anyL{\rho_1'}"] & \\
        {\anyL{\sigma_2}} \ar[rr, "{\ext^{\anyL{\rho_2}}\anyL{\rho_1}}"'] \ar[rrrd, bend right=20, "\anyL{\rho_2'}"'] 
            && {\anyL{\sigma_1} \liplus{\anyL\sigma} \anyL{\sigma_2}} \ar[dr, dashed, "\anyL{\hat\rho}"] & \\
                                                    &&& {\anyL{\hat\sigma}}
    \end{tikzcd}
    \end{center}
    which by \cref{lem:gen-cospan} extends to a unique $\hat\rho$.
    $\mcatI$-neutrality of the injections and composition of independent (product) squares follows from the independent structure on $\worldC$.
\end{proof}

\begin{example}\label{ex:tempC-lic}
The following diagram is a L.I.C. square: 
\begin{center}
    \begin{tikzcd}
                            & {?_{\va,\vb,\vc}} \ar[ld, "\rho_1 = \pair{?_{\va, \vb}}{\incl^{\va,\vb,\vc \inclto \va,\vb,\vc,\vd}}"'] 
                                                \ar[rd, "\rho_2 = \pair{?_{\va,\vc}}{\incl^{\va,\vb,\vc \inclto \va,\vb,\vc,\ve}}"] & \\
    {\incl^{\va, \vb \inclto \va,\vb,\vc,\vd}} \ar[rd, "\ext^{\rho_1}{\rho_2} = \pair{\incl^{\va,\vb \inclto \va,\vb,\vc}}{\incl^{\va,\vb,\vc,\vd \inclto \va,\vb,\vc,\vd,\ve}}"'] && 
    {\incl^{\va,\vc \inclto \va,\vb,\vc,\ve}} \ar[ld, "\ext^{\rho_2}{\rho_1} = \pair{\incl^{\va,\vc \inclto \va,\vb,\vc}}{\incl^{\va,\vb,\vc,\ve \inclto \va,\vb,\vc,\vd,\ve}}"] \\
                                                                            & {\incl^{\va,\vb,\vc \inclto \va,\vb,\vc,\vd,\ve}}
    \end{tikzcd}
\end{center}
Notice that location variables $\vd$ and $\ve$ are not conflated because they are fresh w.r.t. $?_{\va,\vb,\vc}$.
\end{example}

\begin{remark}\label{rem:lic-T-data}
Applying \cref{lem:gen-cospan-data} to the L.I.C. of morphisms $\sigma_1 \xfrom{\rho_1} \sigma \xto{\rho_2} \sigma_2$,
we have that if $\rho_1$ is equipped with instantiation data $\mvarM : \dheapF(\concL{\sigma_1} \ominus \concL{\rho_1},\anyL{\sigma_1})$,
then $\ext^{\rho_2}\rho_1$ can be equipped with instantiation data 
\[ \mvarM_{\ext^{\rho_2}\rho_1} : \dheapF(\concL{\sigma_1 \liplus\sigma \sigma_2} \ominus \concL{\ext^{\rho_2}\rho_1}, \anyL{\sigma_1 \liplus\sigma \sigma_2}) \] 
\end{remark}

\begin{example}
    Consider the square from \cref{ex:tempC-lic}.
    If morphism $\rho_1$ is equipped with instantiation data $\mvarM : \dheapF(\{\va, \vb\}, \{\va,\vb,\vc,\vd\})$,
    we equip $\ext^{\rho_2}\rho_1$ with instantiation data $\dheapF(\{\vb \}, \{\va,\vb,\vc, \vd, \ve\})$ by $\dheapF(\incl^{\vb \inclto \va,\vb}, \incl^{\va,\vb,\vc,\vd \inclto \va,\vb,\vc,\vd,\ve})(\mvarM)$.
\end{example}

\begin{lemma}[Pasting Law for Local Independent Coproducts]\label{lem:lic-pasting}
    In a diagram
    \begin{tikzcd}
        {\sigma} \ar[r, "\rho_1"] \ar[d, "\rho_1'"']
            & {\sigma_1} \ar[r, "\rho_3"] \ar[d, "\rho_2"]
                & {\sigma_3} \ar[d, "\rho_3'"]
        \\
        {\sigma_1'} \ar[r, "\rho_2'"']
            & {\sigma_2} \ar[r, "\rho_3''"']
                & {\sigma_3'}
    \end{tikzcd}
    where the left square is a L.I.C., then the whole square is a L.I.C. iff the right square is.
\end{lemma}
\begin{proof}
    By \cref{lem:gen-cospan-epi-prop}, to show that a commuting square is a L.I.C., it is enough to show that the morphisms out of the span are jointly surjective at the $\anyL\dash$, and $\concL\dash$ level, and further that we have a pullback at the $\anyL\dash$ level.

    For the $\to$ direction, it is clear that if any $\anyL{\rho_3''}$ and $\anyL{\rho_3'}$ are not jointly surjective then the whole square cannot be as well which is a contradiction.
    The same applies to the $\concL\dash$ level, and so the right square must be jointly surjective.
    We must further show that we have a pullback at the $\anyL\dash$ level, and so consider $\anyL{\rho_3''}(\vb) = \anyL{\rho_3'}(\va)$.
    From the left square being a L.I.C. we have that $\anyL{\rho_2'}$ and $\anyL{\rho_2}$ are jointly surjective.
    So if $\anyL{\rho_2}(\vb') = \vb$ then we are done. 
    Otherwise we have $\anyL{\rho_2'}(\vb') = \vb$, 
    and by the outer square being L.I.C., it is a pullback and therefor, 
    $\vb'$ and $\va$ are related by $\anyL{\sigma}$, which implies that they are related by $\anyL{\sigma_1}$ also.
    \nl
    For the $\from$ direction, consider a $\va : \vC \in \anyL{\sigma_3'}$, and we show that it is mapped from $\anyL{\sigma_1'}$ or $\anyL{\sigma_3}$.
    By the right square being jointly surjective, either $\va$ is mapped from $\anyL{\sigma_3}$ or it comes from $\anyL{\sigma_2}$.
    If it comes from $\anyL{\sigma_2}$, then by the left square being L.I.C. either it is mapped from $\anyL{\sigma_1'}$ or it is mapped from $\anyL{\sigma_1}$ and since it is a commuting square, it must also be mapped from $\anyL{\sigma_3}$.
    The same reasoning applies to $\concL\dash$.
    That the outer square is also a pullback at the $\anyL\dash$ level follows from pullback pasting.
\end{proof}

\subsubsection{Variable Preserving Morphisms}

We identify a class of morphisms in $\tempC$: morphisms that do not instantiate any location variables from $\varL{\sigma_1}$.
We call these morphisms variable preserving since every variable in $\varL{\sigma_1}$ is also a variable in $\varL{\sigma_2}$.
\begin{definition}[Variable Preserving Morphisms]
A morphism $\rho : \sigma_1 \to \sigma_2$ is variable preserving if it satisfies $\concL{\sigma_2} \pltimes{\anyL{\sigma_2}} \varL{\sigma_1} = \emptyset$.
\end{definition}
\begin{example}
    The morphism $?_{\va} \xto{\pair{?_{\vb}}{\incl^{\va \inclto \va,\vb,\vc}}} \incl^{\vb \inclto \va,\vb,\vc}$ is a variable preserving morphism since variable $\va$ from the domain is mapped to a variable in the codomain.
    The morphism $?_\va \xto{\pair{?_\va}{1_\va}} 1_\va$ is not a variable preserving morphism since variable $\va$ is instantiated by the morphism.
\end{example}

\begin{lemma}[Variable Preserving Property Equivalence]\label{lem:vpres-prop}
    Let $\rho : \sigma_1 \to \sigma_2$. We have the following equivalences

    \begin{tikzcd}
        {\emptyset} 
            \ar[r] \ar[d] \pbsquare{dr} 
            & {\varL{\sigma_1}} 
                \ar[d, "\varL\rho"'] 
        \\
        {\concL{\sigma_2}} 
            \ar[r, "\concL{\sigma_2}"'] 
            & {\anyL{\sigma_2}}
    \end{tikzcd} $\xLeftrightarrow{(1)}$
    \begin{tikzcd}
        {\anyL{\sigma_1}} 
            \ar[r, "\anyL\rho"]
            & {\anyL{\sigma_2}}
        \\
        {\concL{\sigma_1}} 
            \ar[u, "\concL{\sigma_1}"] \ar[r, "\concL\rho"] \pbsquare{ur}
            & {\concL{\sigma_2}} 
                \ar[u, "\concL{\sigma_2}"]
    \end{tikzcd} $\xLeftrightarrow{(2)}$
    \begin{tikzcd}
        {\varL{\sigma_1}} 
            \ar[dr, "\varL\rho"'] \ar[r, dashed, "\exists \varL\rho\res|{\sigma_2}"] 
            & {\varL{\sigma_2}} 
                \ar[d, "\varL{\sigma_2}"] 
        \\
            & {\anyL{\sigma_2}}
    \end{tikzcd} $\xLeftrightarrow{(3)}$
    \begin{tikzcd}
        {\concL{\sigma_2} \ominus \concL\rho} 
            \ar[d, dashed, "\exists \concL{\sigma_2}\res|\rho"] \ar[r, "\compl{\concL\rho}"] 
            & {\concL{\sigma_2}} 
                \ar[d, "\concL{\sigma_2}"'] 
        \\
        {\anyL{\sigma_2} \ominus \anyL\rho} 
            \ar[r, "\compl{\anyL\rho}"'] 
            & {\anyL{\sigma_2}}
    \end{tikzcd}

\end{lemma}
\begin{proof}\nl
    For (1), consider the canonical isomorphism $\anyL{\sigma_1} \iso \concL{\sigma_1} \mplus \varL{\sigma_1}$.

    For $\xRightarrow{(2)}$ consider $\varL\rho : \varL{\sigma_1} \to \anyL{\sigma_2} \iso \concL{\sigma_2} \mplus \varL{\sigma_2}$.
    If $\varL\rho$ maps to $\concL{\sigma_2}$ we get a contradiction from the pullback.

    For $\xLeftarrow{(2)}$, we have the diagram 
    \begin{tikzcd}
        {\emptyset} \ar[d] \ar[r] & {\varL{\sigma_1}} \ar[d, "\varL\rho"'] \ar[r, dashed, "\varL\rho\res|{\sigma_2}"] & {\varL{\sigma_2}} \ar[dl, "\varL{\sigma_2}"] \\
        {\concL{\sigma_2}} \ar[r, "\concL{\sigma_2}"] & {\anyL{\sigma_2}}
    \end{tikzcd}
    which is a pullback since $\img{\varL{\sigma_2}} \setInt \img{\concL{\sigma_2}} = \emptyset$ which makes the inner square a pullback.

    For $\xRightarrow{(3)}$ the second pullback gives $\concL{\sigma_2}\compl{\concL\rho}(\va : \vC) \not\in \img{\anyL\rho}$.

    For $\xLeftarrow{(3)}$, we show $\concL{\sigma_2} \pltimes{\anyL{\sigma_2}} \varL{\sigma_1} = \emptyset$, and so consider $\varL\rho(\va : \vC) = \concL{\sigma_2}(\vb : \vC)$.
    From $\rho$ being a commuting square we have $\vb : \vC \in \concL{\sigma_2} \ominus \concL\rho$. 
    But then $\concL{\sigma_2}(\vb) \in \anyL{\sigma_2} \ominus \anyL\rho$ which is a contradiction as it must also be equal to $\varL\rho(\va)$.
\end{proof}

\begin{notation}
    For a variable preserving morphism $\sigma_1 \xto\rho \sigma_2$, 
    let $\compl\rho : (\sigma_2 \ominus \rho) \to \sigma_2$ denote the morphism in $\tempC$ given by $\xRightarrow{(3)}$ in \cref{lem:vpres-prop}.
\end{notation}

\begin{lemma}[Variable Preservation Closed Under Composition]\label{lem:comp-vpres}
    Variable preserving morphisms are closed under composition.
\end{lemma}
\begin{proof}
    Follows directly by pullback pasting.
\end{proof}

\begin{lemma}[Variable Preserving Cancellation]\label{lem:vpres-cancl}
    In \begin{tikzcd}
        {\sigma_1} 
            \ar[rr, "\rho"] \ar[dr, "\rho_1"']
                && {\sigma_3}
        \\
            & {\sigma_2}
                \ar[ur, "\rho_2"']
    \end{tikzcd}
    if $\rho$ and $\rho_2$ are variable preserving, so is $\rho_1$.
\end{lemma}
\begin{proof}
    Follows directly by pullback pasting.
\end{proof}

\begin{lemma}[Equalizers are Variable Preserving]\label{lem:eq-vpres-tempC}
   In an equalizer diagram
    \begin{center}
    \begin{tikzcd}
        {\sigma_1} \ar[r, "\ve"] & {\sigma_2} \ar[rr, shift left=1,"\rho_1"] \ar[rr, shift right=1, "\rho_2"'] && {\sigma_3}
    \end{tikzcd}
    \end{center}
    $\ve$ is variable preserving.
\end{lemma}
\begin{proof}\nl
    We want to show $\concL{\sigma_2} \pltimes{\anyL{\sigma_2}} \varL{\sigma_1} = \emptyset$ and so consider
    \begin{center}
        \begin{tikzcd}
            {\vw} \ar[rrd, bend left=30, "\tau_2"] \ar[ddr, bend right=30, "\tau_1"'] && \\
                                                          & {\emptyset} \ar[r] \ar[d]                      & {\varL{\sigma_1}} \ar[d, "\varL\ve"] \\
                                                          & {\concL{\sigma_2}} \ar[r, "\concL{\sigma_2}"'] & {\anyL{\sigma_2}}
        \end{tikzcd}
    \end{center}
    Notice that the following square $\ve'$ commutes
    \begin{center}
        \begin{tikzcd}
            {\concL{\sigma_1} \mplus \varL{\sigma_1} \iso \anyL{\sigma_1}} \ar[r, "\anyL\ve"] & {\anyL{\sigma_2}} \\
            {\concL{\sigma_1} \mplus \vw} \ar[u, "1 \mplus \tau_2"] \ar[r, "\copair{\concL\ve}{\tau_1}"'] & {\concL{\sigma_2}} \ar[u, "\concL{\sigma_2}"']
        \end{tikzcd}
    \end{center}
    and the copairing $\copair{\concL\ve}{\tau_1}$ is mono
    \begin{align*}
        \copair{\concL\ve}{\tau_1} \comp \theta_1 & = \copair{\concL\ve}{\tau_1} \comp \theta_2 \iff \\
        \concL{\sigma_2} \comp \copair{\concL\ve}{\tau_1} \comp \theta_1 & = \concL{\sigma_2} \comp \copair{\concL\ve}{\tau_1} \comp \theta_2 \iff \\
        \anyL\ve \comp 1 \oplus \tau_2 \comp \theta_1 & = \anyL\ve \comp 1 \oplus \tau_2 \comp \theta_2 \iff \\
        1 \oplus \tau_2 \comp \theta_1 & = 1 \oplus \tau_2 \comp \theta_2 \iff \\
        \theta_1 & = \theta_2
    \end{align*}

    Further notice that $\ve'$ commutes with $\rho_1$, $\rho_2$ from $\anyL{\rho_1} \comp \anyL\ve = \anyL{\rho_1} \comp \anyL\ve$
    \begin{align*}
        \concL{\rho_1} \comp \tau_1 &= \concL{\rho_2} \comp \tau_1 \iff \\
        \concL{\sigma_3} \comp \concL{\rho_1} \comp \tau_1   &= \concL{\sigma_3} \comp \concL{\rho_2} \comp \tau_1 \iff \\
        \anyL{\rho_1} \comp \concL{\sigma_2} \comp \tau_1      &= \anyL{\rho_2} \comp \concL{\sigma_2} \comp \tau_1 \iff \\
        \anyL{\rho_1} \comp \varL\ve \comp \tau_2  &= \anyL{\rho_2} \comp \varL\ve \comp \tau_2 =  \\
        \anyL{\rho_1} \comp \anyL\ve \comp \varL{\sigma_1} \comp \tau_2  &= \anyL{\rho_2} \comp \anyL\ve \comp \varL{\sigma_1} \comp \tau_2
    \end{align*}
    and so the universal property of the equalizer gives us an injective morphism $\concL{\sigma_1} \mplus \vw \to \concL{\sigma_1}$, which means $\vw = \emptyset$. 
\end{proof}

\begin{lemma}[Regular Monos in $\tempC$]\label{lem:rmono-vpres-tempC}
   A morphism $\rho : \sigma_1 \to \sigma_2$ in $\tempC$ is regular mono iff $\rho$ is variable preserving. 
\end{lemma}
\begin{proof}
    \nl
    $\to$) Follows directly from \cref{lem:eq-vpres-tempC}.
    \nl
    $\from$)
    We show that this is an equalizer diagram: 
    \begin{center}
    \begin{tikzcd}
        {\sigma_1} \ar[r, "{\rho}"] & {\sigma_2} \ar[rr, shift left=1,"{\ext^\rho\rho}_l"] \ar[rr, shift right=1, "{\ext^\rho\rho}_r"'] && {\sigma_2 \liplus{\sigma_1} \sigma_2}
    \end{tikzcd}
    \end{center}
    
    From \cref{lem:vpres-prop} and \cref{lem:lic-vpres} we have $\concL{\sigma_2}\res|\rho$ and we can simplify the diagram to
    \begin{center}
    \begin{tikzcd}
            & {\anyL{\sigma_1} \mplus (\anyL{\sigma_2} \ominus \anyL\rho)} 
            \ar[rr, shift left=1, "1 \mplus \incl_1"] 
            \ar[rr, shift right=1, "1 \mplus \incl_2"'] 
                && {\anyL{\sigma_1} \mplus (\anyL{\sigma_2} \ominus \anyL\rho) \mplus (\anyL{\sigma_2} \ominus \anyL\rho)} & 
\\
    {\sigma_2} \ar[r, phantom, pos=0.20, "\iso"] & {} && {} \ar[r, phantom, pos=0.95, "\iso"] & {\sigma_2 \liplus{\sigma_1} \sigma_2} 
\\
            & {\concL{\sigma_1} \mplus (\concL{\sigma_2} \ominus \concL\rho)} 
            \ar[uu, "\concL{\sigma_1} \mplus \concL{\sigma_2}\res|\rho"'] 
            \ar[rr, shift left=1, "1 \mplus \incl_1"] 
            \ar[rr, shift right=1, "1 \mplus \incl_2"']    
                                                        && {\concL{\sigma_1} \mplus (\concL{\sigma_2} \ominus \concL\rho) \mplus (\concL{\sigma_2} \ominus \concL\rho)} 
        \ar[uu, "\concL{\sigma_1} \mplus \concL{\sigma_2}\res|\rho \mplus \concL{\sigma_2}\res|\rho"'] & \\
    \end{tikzcd}
    \end{center}
    Consider $(1 \mplus \incl_1) \comp \iso \comp \phi = (1 \mplus \incl_2) \comp \iso \comp \phi$. Then $\phi$ must factor through $\rho$ and uniqueness follows from $\rho$ being mono.
\end{proof}

\begin{lemma}[Isos are Variable Preserving]\label{lem:iso-vpres}
    All isomorphisms are variable preserving morphisms.
\end{lemma}
\begin{proof} Follows from \cref{lem:rmono-vpres-tempC}. \end{proof}

\begin{lemma}[Local Independent Coproducts and Variable Preserving Property]\label{lem:lic-vpres}
    In a L.I.C. where $\rho_2$ is a variable preserving morphism
    \begin{center}
    \begin{tikzcd}[sep=small]     
                                                     & {\sigma} \ar[ld, "{\rho_1}"'] \ar[rd, "{\rho_2}" black]  & \\
        {\sigma_1} \ar[rd, "{\ext^{\rho_1}\rho_2}"'] &                                                                     & {\sigma_2} \ar[ld, "{\ext^{\rho_2}\rho_1}"] \\
                                                     &              {\sigma_1 \liplus\sigma \sigma_2}                      &
    \end{tikzcd}
    \end{center}
    we have that $\sigma_1 \liplus\sigma \sigma_2 \iso \concL{\sigma_1} \liplus{\concL\sigma} \concL{\sigma_2} \xdashto{\concL{\sigma_1} \mplus \concL{\sigma_2}\res|{\rho_2}}
                                                                                                       \anyL{\sigma_1} \liplus{\anyL\sigma} \anyL{\sigma_2}$ and
                 $\ext^{\rho_1}\rho_2$ is variable preserving as well.
\end{lemma}
\begin{proof}\nl
    Notice that this is a pullback square
    \begin{center}
    \begin{tikzcd}[sep=small]
        {\concL\sigma} \ar[rrrr, "\concL{\rho_1}"] \ar[dd, "\concL{\rho_2}"'] \pbsquare{ddrrrr}
                  &&&& {\concL{\sigma_1}} \ar[d, "\concL{\sigma_1}"] \\
                  &&&& {\anyL{\sigma_1}} \ar[d, "\ext^{\anyL{\rho_1}}\anyL{\rho_2}"] \\
        {\concL{\sigma_2}} \ar[r, "\concL{\sigma_2}"'] 
                  & {\anyL{\sigma_2}} \ar[rrr, "\ext^{\anyL{\rho_2}}\anyL{\rho_1}"']
                  &&& {\anyL{\sigma_1} \liplus{\anyL\sigma} \anyL{\sigma_2}}
    \end{tikzcd}
    \end{center}
    By noticing the canonical isomorphism 
    \[ \anyL{\sigma_1} \liplus{\anyL\sigma} \anyL{\sigma_2} \iso \concL{\sigma} \oplus \varL\sigma \oplus (\anyL{\sigma_1} \ominus \anyL{\rho_1})
                                                                                               \oplus (\anyL{\sigma_2} \ominus \anyL{\rho_2}) \]
    and that $\ext^{\anyL{\rho_1}}\anyL{\rho_2} \comp \concL{\sigma_1}$ maps to $\concL\sigma$, $\varL\sigma$ and $\anyL{\sigma_1} \ominus \anyL{\rho_1}$,
    whereas $\ext^{\anyL{\rho_2}}\anyL{\rho_1} \comp \concL{\sigma_2}$ maps to $\concL\sigma$ and $(\anyL{\sigma_2} \ominus \anyL{\rho_2})$ only (via $\concL{\sigma_2}\res|{\rho_2}$).

    That the induced morphism $\concL{\sigma_1} \liplus{\concL{\sigma}} \concL{\sigma_2}$ is $\concL{\sigma_1} \mplus \concL{\sigma_2}\res|{\rho_2}$ follows from uniqueness 
    ($\concL{\sigma_1}$ here has the isomorphic type $\concL\sigma \mplus (\concL{\sigma_1} \ominus \concL{\rho_1}) \to \anyL\sigma \mplus (\anyL{\sigma_1} \ominus \anyL{\rho_1})$). 

    That $\ext^{\rho_1}\rho_2$ is variable preserving follows from
    \begin{center}
    \begin{tikzcd}
        {(\concL{\sigma_1} \liplus{\concL\sigma} \concL{\sigma_2}) \ominus \ext^{\concL{\rho_1}}\concL{\rho_2}} 
            \ar[dd, bend right=40, "\iso"', end anchor=west] \ar[r, "\compl{\ext^{\concL{\rho_1}}\concL{\rho_2}}" yshift=5pt]
            & {\concL{\sigma_1} \liplus{\concL\sigma} \concL{\sigma_2}} 
                \ar[r, "\concL{\sigma_1} \liplus{\concL\sigma} \concL{\sigma_2}"]
                & {\anyL{\sigma_1} \liplus{\anyL\sigma} \anyL{\sigma_2}} 
        \\
        {\concL{\sigma_2}} 
            \ar[r, "\concL{\sigma_2}"] \ar[ur, "\ext^{\concL{\sigma_2}}\concL{\sigma_1}"]
            & {\anyL{\sigma_2}} 
                \ar[ru, "\ext^{\anyL{\rho_2}}\anyL{\rho_1}"']
        \\
        {\concL{\sigma_2} \ominus \concL{\rho_2}} 
            \ar[u, "\compl{\concL{\rho_2}}"'] \ar[r, "\concL{\sigma_2}\res|{\rho_2}"'] 
            & {\anyL{\sigma_2} \ominus \anyL{\rho_2}} 
                \ar[u, "\compl{\anyL{\rho_2}}"'] \ar[r, "\iso"']
                & {(\anyL{\sigma_1} \liplus{\anyL\sigma} \anyL{\sigma_2}) \ominus \ext^{\anyL{\rho_1}}\anyL{\rho_2}} 
                    \ar[uu, "\compl{\ext^{\anyL{\rho_1}}\anyL{\rho_2}}"']
    \end{tikzcd}
    \end{center}
\end{proof}

\begin{lemma}[Pullbacks and Variable Preserving Property]\label{lem:pullback-vpres}
    Let $\rho_2$ be variable preserving in 
    \begin{center}
    \begin{tikzcd}[sep=small]     
            & {\sigma} \ar[ld, "\rho_1"'] \ar[rd, "\rho_2"] 
        \\
        {\sigma_1} \ar[rd, "\rho_1'"'] 
                && {\sigma_2} \ar[ld, "\rho_2'"] 
        \\
            & {\hat\sigma}
    \end{tikzcd}
    \end{center}
    Then the above diagram is a pullback iff its $\anyL\dash$-image is a pullback.
\end{lemma}
\begin{proof}
    The $\to$ direction is immediate by \cref{lem:pullbacks-tempC}.
    For $\from$, by \cref{lem:pullback-mono-post-comp} it is enough to show that
    \begin{tikzcd}[sep=small]
        {\concL\sigma} 
            \ar[rr, "\concL{\rho_1}"] \ar[dd, "\concL{\rho_2}"]
            && {\concL{\sigma_1}}
                \ar[d, "\concL{\rho_1'}"]
        \\
            && {\concL{\hat\sigma}}
                \ar[d, "\concL{\hat\sigma}"]
        \\
        {\concL{\sigma_2}}
            \ar[r, "\concL{\rho_2'}"]
                & {\concL{\hat\sigma}}
                    \ar[r, "\concL{\hat\sigma}"]
                    & {\anyL{\hat\sigma}}
    \end{tikzcd}
    is a pullback which is the same as
    \begin{tikzcd}[sep=small]
        {\concL\sigma} 
            \ar[rr, "\concL{\rho_1}"] \ar[dd, "\concL{\rho_2}"]
            && {\concL{\sigma_1}}
                \ar[d, "\concL{\sigma_1}"]
        \\
            && {\anyL{\sigma_1}}
                \ar[d, "\anyL{\rho_1'}"]
        \\
        {\concL{\sigma_2}}
            \ar[r, "\concL{\sigma_2}"]
                & {\anyL{\sigma_2}}
                    \ar[r, "\anyL{\rho_2'}"]
                    & {\anyL{\hat\sigma}}
    \end{tikzcd}
    which is given by the pullback
    \begin{tikzcd}
        {\concL\sigma} 
            \ar[r, "\concL\sigma"] \ar[d, "\concL{\rho_2}"'] \pbsquare{dr}
            & {\anyL\sigma}
                \ar[r, "\anyL{\rho_1'}"] \ar[d, "\anyL{\rho_2}"] \pbsquare{dr}
                    & {\anyL{\sigma_1}}
                        \ar[d, "\anyL{\rho_1'}"]
        \\
        {\concL{\sigma_2}}
            \ar[r, "\concL{\sigma_2}"] 
                & {\anyL{\sigma_2}}
                    \ar[r, "\anyL{\rho_2'}"]
                    & {\anyL{\hat\sigma}}
    \end{tikzcd} and \cref{lem:pullback-factoring}.
\end{proof}

\subsubsection{Variable Reflecting Morphisms}

We identify another class of morphisms in $\tempC$: morphisms that do not introduce location variables.
We call these morphisms variable reflecting since every variable in $\varL{\sigma_2}$ comes from $\varL{\sigma_1}$.
\begin{definition}[Variable Reflecting Morphisms]
A morphism $\rho : \sigma_1 \to \sigma_2$ is variable reflecting if it satisfies $(\anyL{\sigma_2} \ominus \varL\rho) \pltimes{\anyL{\sigma_2}} \varL{\sigma_2} = \emptyset$.
\end{definition}
\begin{example}
    The morphism $?_{\va,\vb} \xto{\pair{?_\va}{1_{\va,\vb}}} \incl^{\va \inclto \va,\vb}$ is a variable reflecting morphism since variable $\vb$ from the codomain comes from the domain.
    The morphism $?_\va \xto{\pair{?_\va}{\incl^{\va \inclto \va,\vb}}} \incl^{\va \inclto \va, \vb}$ is not a variable reflecting morphism since variable $\vb$ is fresh.
\end{example}

\begin{lemma}[Variable Reflecting Property Equivalence]\label{lem:vrefl-prop}
    Let $\rho : \sigma_1 \to \sigma_2$. We have the following equivalences
    \begin{center}
    \begin{tikzcd}
        {\emptyset} \ar[r] \ar[d] \pbsquare{dr} & {\anyL{\sigma_2} \ominus \varL\rho} \ar[d, "\compl{\varL\rho}"'] \\
        {\varL{\sigma_2}} \ar[r, "\varL{\sigma_2}"'] & {\anyL{\sigma_2}}
    \end{tikzcd} $\xLeftrightarrow{(1)}$
    \begin{tikzcd}
        {\varL{\sigma_2}} \ar[rr, dashed, "\exists \inv^{\varL\rho}\res|{\sigma_2}"] \ar[drr, "\varL{\sigma_2}"'] && {\varL{\sigma_1}} \ar[d, "\varL\rho"'] \\
                                                                                                                   && {\anyL{\sigma_2}}
    \end{tikzcd} $\xLeftrightarrow{(2)}$
    \begin{tikzcd}
        {\concL{\sigma_2} \ominus \concL\rho} \ar[r, "\compl{\concL\rho}"] & {\concL{\sigma_2}} \ar[d, "\concL{\sigma_2}"] \\
        {\anyL{\sigma_2} \ominus \anyL\rho} \ar[u, dashed, "\exists \inv^{\concL{\sigma_2}}\res|\rho"] \ar[r, "\compl{\anyL\rho}"'] & {\anyL{\sigma_2}}
    \end{tikzcd}
    \end{center}
\end{lemma}
\begin{proof}
    For $\xRightarrow{(1)}$, consider $\varL\rho : \varL{\sigma_1} \to \anyL{\sigma_2} \iso \concL{\sigma_2} \mplus \varL{\sigma_2}$.
    If there exists $\va : \vC \in \varL{\sigma_2}$ that $\varL\rho$ does not map to, it must also be an element of $\anyL{\sigma_2} \ominus \varL\rho$ by definition. 
    But then the pullback gives us a contradiction.
    For $\xLeftarrow{(1)}$, we have a diagram
    \begin{center}
    \begin{tikzcd}
        {\emptyset} \ar[d] \ar[r] & {\anyL{\sigma_2} \ominus \varL\rho} \ar[d, "\compl{\varL\rho}"] \\
        {\varL{\sigma_2}} \ar[r, "\varL{\sigma_2}"] \ar[dr, "\inv^{\varL\rho}\res|{\sigma_2}"'] & {\anyL{\sigma_2}} \\
                                                                                                & {\varL{\sigma_1}} \ar[u, "\varL\rho"']
    \end{tikzcd}
    \end{center}
    which is a pullback since $\img{\varL\rho} \setInt \img{\compl{\varL\rho}} = \emptyset$.
    For $\xRightarrow{(2)}$, first notice that this is a pullback
    \begin{center}
    \begin{tikzcd}
        {\emptyset} \ar[d] \ar[r] \pbsquare{dr} & {\anyL{\sigma_2} \ominus \varL\rho} \ar[d, "\compl{\anyL\rho}"] \\
        {\varL{\sigma_2}} \ar[r, "\varL{\sigma_2}"'] & {\anyL{\sigma_2}}
    \end{tikzcd}
    \end{center}
    given by the inclusion $\anyL{\sigma_2} \ominus \anyL\rho \inclto \anyL{\sigma_2} \ominus \varL\rho$ and pullback pasting.
    Now consider the morphism
    \[ \anyL{\sigma_2} \ominus \anyL\rho \xto{\compl{\anyL\rho}} \anyL{\sigma_2} \iso \concL{\sigma_1} \mplus (\concL{\sigma_2} \ominus \concL\rho) \mplus \varL{\sigma_2} \]
    From $\rho$ being a commuting square it cannot map to $\concL{\sigma_1}$, and from the pullback it can also not map to $\varL{\sigma_2}$.
    Lastly for $\xLeftarrow{(2)}$, consider $\varL\rho : \varL{\sigma_1} \to \anyL{\sigma_2} \iso \concL{\sigma_2} \mplus \varL{\sigma_2}$, and a $\va : \vC \in \varL{\sigma_2}$ that is not mapped to from $\varL\rho$.
    This gives us that either $\va \in \anyL{\sigma_2} \ominus \anyL\rho$ which is a contradiction from $\inv^{\concL{\sigma_2}}\res|\rho$
    or it comes from $\concL{\sigma_1}$ which is a contradiction from $\rho$ being a commutative square.

\end{proof}

\begin{lemma}[Composition of Variable Reflecting Morphisms]\label{lem:comp-vrefl}
    Variable reflecting morphisms are closed under composition.
\end{lemma}
\begin{proof}
    For variable reflecting morphisms $\sigma_1 \xto{\rho_1} \sigma_2 \xto{\rho_2} \sigma_3$ we have
    \begin{center}
    \begin{tikzcd}
    {\varL{\sigma_3}} \ar[rd, "\varL{\sigma_3}"'] \ar[r, dashed, "\inv^{\varL{\rho_2}}\res|{\sigma_3}"] 
        & {\varL{\sigma_2}} \ar[d, "\varL{\rho_2}"'] \ar[rd, "\varL{\sigma_2}"'] \ar[r, dashed, "\inv^{\varL{\rho_1}}\res|{\sigma_2}"] 
            & {\varL{\sigma_1}} \ar[d, "\varL{\rho_1}"] \ar[dd, bend left=70, "\varL{\rho_2 \comp \rho_1}"]
    \\
        & {\anyL{\sigma_3}} \ar[rd, Rightarrow]
            & {\anyL{\sigma_2}} \ar[d, "\anyL{\rho_2}"]
    \\
           && {\anyL{\sigma_3}}
    \end{tikzcd}
    \end{center}\vspace{-1.5\baselineskip}
\end{proof}

\begin{lemma}[Variable Reflecting Cancellation]\label{lem:vrefl-cancl}
    In \begin{tikzcd}
        {\sigma_1} 
            \ar[rr, "\rho"] \ar[dr, "\rho_1"']
                && {\sigma_3}
        \\
            & {\sigma_2}
                \ar[ur, "\rho_2"']
    \end{tikzcd}
    if $\rho$ is variable reflecting, so is $\rho_2$.
\end{lemma}
\begin{proof}
    We show
    \begin{tikzcd}
        {\concL{\sigma_3} \ominus \concL{\rho_2}} \ar[r, "\compl{\concL{\rho_2}}"] & {\concL{\sigma_3}} \ar[d, "\concL{\sigma_3}"] \\
        {\anyL{\sigma_3} \ominus \anyL{\rho_2}} \ar[u, dashed, "\exists \inv^{\concL{\sigma_3}}\res|{\rho_2}"] \ar[r, "\compl{\anyL{\rho_2}}"'] & {\anyL{\sigma_3}}
    \end{tikzcd}.
    The variable reflecting property of $\rho$ gives us a $(\anyL{\sigma_3} \ominus \anyL\rho) \dashto (\concL{\sigma_3} \ominus \concL\rho)$
    and using the canonical isos we have
    \[ \anyL{\sigma_3} \ominus \anyL{\rho_2} \to 
       \anyL{\sigma_3} \ominus \anyL\rho     \to 
       \concL{\sigma_3} \ominus \concL\rho   \iso 
       (\concL{\sigma_3} \ominus \concL{\rho_2}) \oplus (\concL{\sigma_2} \ominus \concL{\rho_1}) \] 
    and its image must lie in $\concL{\sigma_3} \ominus \concL{\rho_2}$ by the commuting square of $\rho_2$.
\end{proof}

\begin{lemma}[Variable Preserving with Reflecting Cancellation]\label{lem:vpres-vrefl-cancl}
    In \begin{tikzcd}
        {\sigma_1} 
            \ar[rr, "\rho"] \ar[dr, "\rho_1"']
                && {\sigma_3}
        \\
            & {\sigma_2}
                \ar[ur, "\rho_2"']
    \end{tikzcd}
    if $\rho$ is variable preserving and $\rho_1$ is reflecting, then $\rho_2$ is variable preserving.
\end{lemma}
\begin{proof}
The variable preserving property of $\rho$ gives us a $\concL{\sigma_3} \ominus \concL\rho \dashto \anyL{\sigma_3} \ominus \anyL\rho$,
the variable reflecting property of $\rho_1$ gives us a $\anyL{\sigma_2} \ominus \anyL{\rho_1} \dashto \concL{\sigma_2} \ominus \concL{\rho_1}$,
and together with the canonical isos we have
\[ 
    \concL{\sigma_3} \ominus \concL{\rho_2} \to
        \concL{\sigma_3} \ominus \concL\rho \to
            \anyL{\sigma_3} \ominus \anyL\rho \iso
                (\anyL{\sigma_3} \ominus \anyL{\rho_2}) \oplus (\anyL{\sigma_2} \ominus \anyL{\rho_1}) \to
                    (\anyL{\sigma_3} \ominus \anyL{\rho_2}) \oplus (\concL{\sigma_2} \ominus \concL{\rho_1})
\]
and its image must lie in $\anyL{\sigma_3} \ominus \anyL{\rho_2}$.
\end{proof}

\begin{lemma}[Variable Reflecting with Preserving Cancellation]\label{lem:vrefl-vpres-cancl}
    In \begin{tikzcd}
        {\sigma_1} 
            \ar[rr, "\rho"] \ar[dr, "\rho_1"']
                && {\sigma_3}
        \\
            & {\sigma_2}
                \ar[ur, "\rho_2"']
    \end{tikzcd}
    if $\rho$ is variable reflecting and $\rho_2$ is variable preserving, then $\rho_1$ is variable reflecting.
\end{lemma}
\begin{proof}
    The variable reflecting property of $\rho$ gives us a $(\anyL{\sigma_3} \ominus \anyL\rho) \dashto (\concL{\sigma_3} \ominus \concL\rho)$,
    and the variable preserving property of $\rho_2$ gives us a $\concL{\sigma_3} \ominus \concL{\rho_2} \dashto \anyL{\sigma_3} \ominus \anyL{\rho_2}$,
    and together with the canonical isos we have
    \[ 
        \anyL{\sigma_2} \ominus \anyL{\rho_1} \to 
            \anyL{\sigma_3} \ominus \anyL\rho \to 
                \concL{\sigma_3} \ominus \concL\rho \iso
                    (\concL{\sigma_3} \ominus \concL{\rho_2}) \oplus (\concL{\sigma_2} \ominus \concL{\rho_1}) \to
                        (\anyL{\sigma_3} \ominus \anyL{\rho_2}) \oplus (\concL{\sigma_2} \ominus \concL{\rho_1})
    \]
    and its image must lie in $\concL{\sigma_2} \ominus \concL{\rho_1}$.
\end{proof}

\begin{lemma}[Local Independent Coproducts and Variable Reflecting Morphisms]\label{lem:lic-vrefl}
    In a L.I.C. square where $\rho_2$ is variable reflecting: 
    \begin{center}
    \begin{tikzcd}[sep=small]     
                                                     & {\sigma} \ar[ld, "{\rho_1}"'] \ar[rd, "{\rho_2}" black]  & \\
        {\sigma_1} \ar[rd, "{\ext^{\rho_1}\rho_2}"'] &                                                          & {\sigma_2} \ar[ld, "{\ext^{\rho_2}\rho_1}"] \\
                                                     &              {\sigma_1 \liplus\sigma \sigma_2}               &
    \end{tikzcd}
    \end{center}
    we have that $\ext^{\rho_1}\rho_2$ is variable reflecting as well.
\end{lemma}
\begin{proof}
    Consider
    \begin{center}
    \begin{tikzcd}[sep=small]
    {\concL\sigma} 
        \ar[rrrrd, bend left=15, "\concL{\rho_2}"] 
        \ar[dddr, bend right=60, "\concL{\rho_1}"'] 
        \ar[rd, dashed, "\vh"]
    \\ 
        & {\concL{\sigma_1} \pltimes{\anyL{\sigma_1} \liplus{\anyL\sigma} \anyL{\sigma_2}} \concL{\sigma_2}} 
            \ar[rrr, "\pi_2"] 
            \ar[dd, "\pi_1"'] 
            \ar[dr, dashed, "\vh'"]
                &&& {\concL{\sigma_2}}
                    \ar[d, "\concL{\sigma_2}"]
    \\
            && {\anyL{\sigma}} 
                    \ar[rr, "\anyL{\rho_2}"] \ar[d, "\anyL{\rho_1}"'] \pbsquare{drr}
            && {\anyL{\sigma_2}} 
                    \ar[d, "\ext^{\anyL{\rho_2}}\anyL{\rho_1}"] 
    \\
        & {\concL{\sigma_1}} 
            \ar[r, "\concL{\sigma_1}"'] 
            & {\anyL{\sigma_1}} 
                \ar[rr, "\ext^{\anyL{\rho_1}}\anyL{\rho_2}"'] 
                && {\anyL{\sigma_1} \liplus{\anyL\sigma} \anyL{\sigma_2}}
    \end{tikzcd}
    \end{center}

    We first show that $\inv^{\concL{\sigma_2}}\res|{\rho_2}$ factors through $\concL{\sigma_2} \ominus \pi_2$
    \begin{center}
    \begin{tikzcd}[sep=small]
    {\anyL{\sigma_2} \ominus \anyL{\rho_2}} 
        \ar[rr, "\inv^{\concL{\sigma_2}}\res|{\rho_2}"] 
        \ar[dr, dashed]
            && {\concL{\sigma_2} \ominus \concL{\rho_2}} 
                \ar[dr, "\compl{\concL{\rho_2}}"]
    \\
        & {\concL{\sigma_2} \ominus \pi_2} 
            \ar[rr, "\compl{\pi_2}"']
                && {\concL{\sigma_2}}
    \end{tikzcd}
    \end{center}
    and so consider $\va : \vC \in \anyL{\sigma_2} \ominus \anyL{\rho_2}$.
    From the diagram we have that $\va : \vC \in \anyL{\sigma_2} \ominus (\concL{\sigma_2} \comp \proj_2)$
    and consider the canonical isomorphism 
    \[\anyL{\sigma_2} \ominus (\concL{\sigma_2} \comp \proj_2) \iso 
        (\anyL{\sigma_2} \ominus \concL{\sigma_2}) \mplus (\concL{\sigma_2} \ominus \proj_2) \]
    Morphism $\inv^{\concL{\sigma_2}}\res|{\rho_2}$ tells us that $\va$ is in the image of $\concL{\sigma_2}$, and so
    it must come from $\concL{\sigma_2} \ominus \proj_2$.

    We then show that $\ext^{\rho_1}\rho_2$ is variable reflecting by observing that the following diagram commutes
    \begin{center}
    \begin{tikzcd}[sep=small]
        {\concL{\sigma_1 \liplus{\sigma} \sigma_2} \ominus \concL{\ext^{\rho_1}\rho_2}}
            \ar[drrr, "\compl{\concL{\ext^{\rho_1}\rho_2}}"]
    \\
        {\concL{\sigma_2} \ominus \pi_2} 
            \ar[u, "\iso"] \ar[r, "\compl{\pi_2}"]
            & {\concL{\sigma_2}} 
                \ar[rr, "\concL{\ext^{\rho_2}\rho_1}"'] \ar[dd, "\concL{\sigma_2}"']
                && {\concL{\sigma_1 \liplus{\sigma} \sigma_2}}
                        \ar[dd, "\concL{\sigma_1 \liplus\sigma \sigma_2}"]
    \\\\
        {\anyL{\sigma_2} \ominus \anyL{\rho_2}} 
            \ar[r, "\compl{\anyL{\rho_2}}"] \ar[uu, dashed]
            & {\anyL{\sigma_2}} 
                \ar[rr, "\anyL{\ext^{\rho_2}\rho_1}"]
                    && {\anyL{\sigma_1 \liplus\sigma \sigma_2}}
    \\
        {\anyL{\sigma_1 \liplus\sigma \sigma_2} \ominus \anyL{\ext^{\rho_1}\rho_2}} 
            \ar[u, "\iso"] \ar[urrr, "\compl{\anyL{\ext^{\rho_1}\rho_2}}"']
    \end{tikzcd}
    \end{center}\vspace{-1.5\baselineskip}
\end{proof}

\subsubsection{Factorisation System}

\begin{lemma}[Epimorphisms in $\tempC$]\label{lem:epis-T}
   A morphism $\rho : \sigma_1 \to \sigma_2$ in $\tempC$ is epi iff $\anyL\rho$ is surjective.
\end{lemma}
\begin{proof}\nl
    $\to$) Consider $\rho : \sigma_1 \to \sigma_2$ such that for every $\vl\,\vr$, $\vl \comp \rho = \vr \comp \rho \implies \vl = \vr$.
           Take the L.I.C. of $\sigma_2 \xfrom\rho \sigma_1 \xto\rho \sigma_2$, which gives us two morphisms ${\ext^\rho\rho}_\vl, \ext^\rho\rho_\vr : \sigma_2 \to \sigma_2 \liplus{\sigma_1} \sigma_2$
           that satisfy ${\ext^\rho\rho}_l \comp \rho = {\ext^\rho\rho}_r \comp \rho$.
           From the epi property of $\rho$, we have that $\anyL{{\ext^\rho\rho}_l} = \anyL{{\ext^\rho\rho}_r}$ which implies that $\anyL\rho$ is surjective by definition of $\anyL{\sigma_1} \liplus{\anyL\sigma} \anyL{\sigma_2}$.
    \nl
    $\from$) Let $\vl, \vr$ such that $\vl \comp \rho = \vr \comp \rho$.
             This means that $\anyL\vl \comp \anyL\rho = \anyL\vr \comp \anyL\rho$ and so $\anyL\vl = \anyL\vr$ which also implies $\concL\vl = \concL\vr$.
\end{proof}

\begin{notation}
    We denote epis using $\eto$ and regular monos with $\rmto$.
\end{notation}

\begin{lemma}[Orthogonal Factorisation System in $\tempC$]
    The classes of epis and regular monos (variable preserving) morphisms gives rise to an orthogonal factorisation system in $\tempC$.
\end{lemma}
\begin{proof}
The intuition here is that the epi part will represent the instantiations of the location variables coming from its domain
and the regular mono part will be the introduction of all other locations.

A morphism $\rho : \sigma_1 \to \sigma_2$ is factored as follows
\begin{center}
\begin{tikzcd}
    {\anyL{\sigma_1}} \ar[r, Rightarrow]
        & {\anyL{\sigma_1}} \ar[r, "\anyL\rho"] 
            & {\anyL{\sigma_2}}  
    \\
    {\concL{\sigma_1}} \ar[u, "\concL{\sigma_1}"] \ar[r,dashed] \ar[rr, bend right=30, "\concL\rho"'] 
        & {\concL{\sigma_2} \pltimes{\anyL{\sigma_2}} \anyL{\sigma_1}} \ar[u] \ar[r] \pbsquare{ur}
            & {\concL{\sigma_2}} \ar[u, "\concL{\sigma_2}"']
\end{tikzcd}
\end{center}
\end{proof}
\begin{example}
    We factor the morphism $?_\va \xto{\pair{?_{\va,\vb}}{\incl^{\va \inclto \va,\vb,\vc}}} \incl^{\va,\vb \inclto \va,\vb,\vc}$ as
    \[ ?_\va \xto{\pair{?_\va}{1_\va}} 1_\va \xto{\pair{\incl^{\va \inclto \va,\vb}}{\incl^{\va \inclto \va,\vb,\vc}}} \incl^{\va,\vb \inclto \va,\vb,\vc} \]
\end{example}
